\section{Conclusion}
DisasterClaw contributes an evaluation component for embodied UAV intelligence: it keeps the geographic scope and reference answer fixed while the agent's action changes the observation. Its episode records connect actions, predicted building evidence, answer calls, and terminal outcomes. The four task families capture a trade-off between local detail and regional coverage that building-classification accuracy alone does not describe.

On the frozen final set, image-aware pipelines outperform the evaluated text-only and permitted-metadata controls. At matched recheck transitions, changed observations yield more corrections than harms, but the ROI-cluster interval for the average difference includes zero. The task-wise point estimates are positive for damage and count and negative for spatial questions. T1\_TASK achieves the highest observed complete-policy accuracy, although its adjusted contrast with the search-enabled A0\_NO\_RECHECK policy also includes zero. In the full MESSI release, existing single-temporal YOLO and SegFormer models were connected to DisasterClaw's evidence and recheck controllers for 409 vegetation regions. The action-enabled arm used 20 lower views and answered 372/409 correctly, compared with 369/409 for paired hold; six answers were corrected and three harmed. On the same 20 triggered cases, an equal-budget repeat-high arm answered 12/20 correctly and the real-low arm answered 15/20, isolating the second-view contribution from the additional answer opportunity. A supplementary random-selection control with the same 20-recheck budget averages 369.20/409 over 20 fixed seeds (range 367--371), below the entropy trigger's 372/409. This supports selective allocation on the evaluated cohort; the seed spread does not establish cross-scene generalization. The 34 IrYamim Test regions had appeared in earlier Qwen-only explorations and are not an untouched holdout; both current arms answer 25/34 correctly. These results show that an adapted controller can gate access to new observations in a real-image offline replay. The net gain is small, and physical flight, vehicle dynamics, and transfer of the unchanged disaster-damage pipeline remain untested.
